% page 260 du fac-simile (image p-259 du scan)
% bloc 162.6 x 233.3 mm ; marge gauche 8.0 mm
\pgc{-12.700mm}{0.000mm}{
\l{\cel{5}252}
\l{}
\l{}
\l{\sou{kadavro}. Korpo mortinta, homala o bestiala. - DEFIS.}
\l{}
\l{\sou{kadencar}. (\sou{trans.}) I.Finar frazo, verso, hemistiko tale, ke}
\l{\cel{3}la voco acentizas ta finalo per plufortigar olua pronun-}
\l{\cel{3}ceso. - II. Produktar ritmo per acentizar simetrala ta}
\l{\cel{3}finali,en muziko od en poezio. - III. Produktar movo ye}
\l{\cel{3}intervali inter-egala ed uniforma, qua imitas la ritmo}
\l{\cel{3}muzikala. - DEFIRS.}
\l{}
\l{\sou{kadeto}. (\sou{armeo).}\cel{3}Singla ek la membri di ta kompanii qui}
\l{\cel{3}konsistis ye la familio-juniori volontaria. - DEFIRS.}
\l{}
\l{\sou{kadmio}. (\sou{kemio}). Korpo simpla, forjebla, duktila, di qua la}
\l{\cel{3}rupto produktas fibri, di qua la koloro e la brilo esas}
\l{\cel{3}olti di stano. -\cel{1}\sou{Simbolo kemiala}\cel{1}: Cd. - DEFIRS.}
\l{}
\l{\sou{kadro.}\cel{2}I. Bordumo irga-forma qua cirkondas pikturo, graburo,}
\l{\cel{3}spegulo. - II. Bordumo menuzura. - III. l. To quo limit-}
\l{\cel{3}izas ceno, temo ; 2.\cel{1}\sou{(armeo}). La korpo ek la oficiri ge-}
\l{\cel{3}nerala, oficiri, sub-oficiri, di la divizioni, secioni,}
\l{\cel{3}e c., imperote da qui la soldati su asemblas. - EFS.}
\l{}
\l{\sou{kaduceo}. Emblemo di Merkurio, vergeto kronizita ye du mikra}
\l{\cel{3}ali, cirkum qua esas du serpenti plektita - simbolo di}
\l{\cel{3}komerco, eloquenteso, paco. - DEFIRS.}
\l{}
\l{\sou{kaduka}.\cel{2}I. Vicina ye sua krulo, dekado. - II. Quan ula defekto}
\l{\cel{3}formala igas nihiligota tre balde. - EFIS.}
\l{}
\l{\sou{kafeo}.\cel{2}(\sou{bot.)}\cel{2}I. Grano ek arboreto (L.\cel{1}\sou{coffeo arabica),}\cel{1}sempre}
\l{\cel{3}verda, ek Arabia, di qua la frukto esas grosa quale cerizo}
\l{\cel{3}e konsistas ye shelo qua envelopas pulpo mucilajatra qua}
\l{\cel{3}ipsa cirkondas du alveoli un-grana ; ta grano, torefaktita,}
\l{\cel{3}muelita ed infuzita, donas drinkajo agreabla, ecitiva e to-}
\l{\cel{3}niziva. - II. Infuzuro ek kafeo-grani torefaktita e mue-}
\l{\cel{3}lita. - DEFIRS.}
\l{}
\l{\sou{kafeino.}\cel{2}Alkaloido naturala, kristaligebla, quan onu extrak-}
\l{\cel{3}tas de kafeo, teo, kakao. - DEFIS.}
\l{}
\l{\sou{kaftano}. Honor-peliso quan la sultani ofris a la oficiri pre-}
\l{\cel{3}cipua, a la ambasadisti stranjera, a personegi tre eminen-}
\l{\cel{3}ta, e c. kom donaco. - DEFIRS.}
\l{}
\l{\sou{kagulo}. Quaza froko sen-manika, qua kovris la kapo e la}
\l{\cel{3}parti cetera di la korpo, e qua provizesis ye aperturi}
\l{\cel{3}qui korespondis a la okuli ed a la boko. - FS.}
\l{}
\l{\sou{kaito}.\cel{2}Puero-ludilo, ek framo lejera qua esas kovrita ye}
\l{\cel{3}papero tensita, quan onu igas acensar aden atmosfero per}
\l{\cel{3}kurar kontre-vento e tenante olu per kordeto.\cel{2}- E.}
}
